After studying splitting fields, we are now in a position to study the symmetries of these fields. In particular, we will be interested in the field automorphisms of splitting fields. As we will see, automorphisms for splitting fields are closely tied to the roots of the polynomial that we are splitting. The collection of such automorphisms will form a group under composition, which will be known as the Galois group of the splitting field.

\begin{definition}[Field Automorphism] \leavevmode \\
    Let $L$ be a field.
    \begin{enumerate}
        \item A field isomorphism $\sigma: K \to K$ is called an automorphism of $K$. The set of all automorphisms of $K$ is denoted by $\Aut{K}$.
        \item We say $\sigma \in \Aut{K}$ fixes an element $\alpha \in K$ if $\sigma(\alpha) = \alpha$. If $F\subseteq K$, then $\sigma$ is said to fix $F$ if $\sigma(f) = f$ for all $f \in F$.
    \end{enumerate}
\end{definition}

\begin{example}
    $\sigma(a) = a ~~\forall a \in K$ is called the trivial automorphism (or the identity automorphism), denoted $\sigma_1$ or $1$.
\end{example}

\begin{note}
    The prime subfield of $K$, call it $F$ (for now) is generated by $1$ and since any automorphism $\tau$, $\tau(1)=1$, it follows that $\tau(a) = a$ for all $a\in F$. Hence any automorphisms of a field $K$ fixes its prime subfield. So, all $\tau \in \Q(\sqrt{2})$ fix $\Q$.
\end{note}

\begin{example}
    $\Aut{\Q} = \set{\sigma_1}$ since $\Q$ is its own prime subfield.
\end{example}

\begin{definition} [$\Aut{K/F}$] \leavevmode \\
    Let $K / F$ be an extension of fields. Let $\Aut{K/F}$ be the collection of automorphisms of $K$ which fix $F$.
\end{definition}

\begin{proposition}
    $\Aut{K}$ form a group with respect to composition and $\Aut{K/F}$ form a subgroup of $\Aut{K}$.
\end{proposition}

\begin{proof}
    Composition is an associative binary operation on bijective morphisms. The identity of $\Aut{K}$ is $\sigma_1$. Isomorphisms have functional inverses, which are isomorphisms. So, $\Aut{K}$ forms a group. Now, let $\sigma, \tau \in \Aut{K/F}$. Then let $a\in F$. We have 
    \begin{align*}
        \sigma \tau^{-1}(a) &= \sigma \tau^{-1}(\tau(a)) \\
        &= \sigma(a) \\
        &= a
    \end{align*}
    Hence, $\Aut{K/F} \leq \Aut{K}.$
    \qed
\end{proof}

\begin{proposition}
    Let $K/F$ be a field extension and let $\alpha \in K$ be algebraic over $F$. Then for any $\sigma \in \Aut{K / F}$, $\sigma(\alpha)$ is a root of the minimal polynomial for $\alpha$ over $F$. I.e., $\Aut{K / F}$ permutes the roots of irreducible polynomials. Equivalently, any polynomial in $F[x]$ which has $\alpha$ as a root will also have $\sigma(\alpha)$ as a root.
\end{proposition}

\begin{proof}
    If $\alpha$ satisfies $\sum_{i=0}^\infty a_i \alpha^i = 0$ where $a_i \in F$, then 
    \begin{align*}
        \sigma\left(\sum_{i=0}^\infty a_i \alpha^i\right) &= \sigma(0) \\
        \implies \sum_{i=0}^\infty \sigma(a_i \alpha^i) &= 0 \\
        \implies \sum_{i=0}^\infty \sigma(a_i) \sigma(\alpha^i) &= 0 \\
        \implies \sum_{i=0}^\infty a_i \sigma(\alpha^i) &= 0
    \end{align*}
    \qed
\end{proof}

\begin{example}
    Consider $K = \Q(\sqrt{2})$. If $\tau \in \Aut{K} = \Aut{K/Q}$, then $\tau(\sqrt{2}) = \pm \sqrt{2}$ since $\pm \sqrt{2}$ are the two roots of $m_{\sqrt{2}, \Q}(x)$. We have two possible automorphisms:
    \begin{align*}
        1: \sqrt{2} &\mapsto \sqrt{2} \\
        \sigma : \sqrt{2} &\mapsto -\sqrt{2}
    \end{align*}
    Note $\sigma^2$:
    \begin{align*}
        \sigma^2(\sqrt{2}) &= \sigma(-\sqrt{2}) \\
        &= -\sigma(\sqrt{2}) \\
        &= -(-\sqrt{2}) \\ 
        &= \sqrt{2} \\
        \implies \sigma^2 &= 1
    \end{align*}
    Then $\Aut{\Q(\sqrt{2})/\Q} = \cyc{\sigma} \cong \quotient{\Z}{2\Z}$.
\end{example}

\begin{example}
    Consider $\Q(\sqrt[3]{2})$. If $\tau \in \aut{K}$, we have the following possibilities: $\tau \in (\sqrt[3]{2}) \in \set{\sqrt[3]{2}, \sqrt[3]{2}\omega, \sqrt[3]{2}\omega^2}$ since $m_{\sqrt[3]{2}, \Q}(x) = x^3 -2.$ Since $\Q(\sqrt[3]{2}) \subseteq \R$, $\tau$ cannot map to  $\sqrt[3]{2}\omega, \sqrt[3]{2}\omega^2 \in \C$. Hence $\tau(\sqrt[3]{2}) = \sqrt[3]{2}$ ($\tau = 1$). Thus there is only one automorphism of $\Q(\sqrt[3]{2})$, namely $1$. Hence, $\aut{\Q(\sqrt[3]{2}) / \Q} = 1.$
\end{example}
Comparing $\aut{\Q(\sqrt{2})}$ and $\aut{\Q(\sqrt[3]{2})}$, we see that $\aut{\Q(\sqrt[3]{2})}$ is missing clear automorphisms. In general, we see that if $K/F$ is generated by some collection of elements, then $\sigma \in \aut{K/F}$ is completely determined by where $\sigma$ maps the generators.

\begin{proposition}
    Let $H \leq \aut{K}$. Then the elements of $K$ fixed by all automorphisms in $H$ is a subfield of $K$.
\end{proposition}

\begin{proof}
    Let $h \in H$, $a,b\in F$. Then, by definition, $h(a) = a$ and $h(b) = b$ so that 
    $$h(a\pm b) = h(a) \pm h(b).$$
    Further,
    $$h(ab) = h(a)h(b)$$
    and
    $$h(a^{-1}) = h(a)^{-1} = a^{-1}$$
    Thus, $F$ is a field.
    \qed
\end{proof}

\begin{note}
    $H$ does not have to be a subgroup for the proposition to hold.
\end{note}

We the above proposition, we take interest in the elements in $K$ that are fixed by the automorphisms of $K$.

\begin{definition}[Fixed Field] \leavevmode \\
    If $H \leq \aut{K}$, the subfield fixed by $H$ is called the fixed field of $H$.
\end{definition}

\begin{proposition}
    The association of groups of $\aut{K}$ to fixed fields and fixed fields to subgroups is inclusion reversing. Namely,
    \begin{enumerate}
        \item If $F_1 \subseteq F_2 \subseteq K$ are subfields of $K$, then $\aut{K/F_2} \subseteq \aut{K /F_1}$.
        \item If $H_1 \leq H_2 \leq \aut{K}$ where $F_1$ and $F_2$ are the associated fixed fields of $H_1$ and $H_2$, respectively, then $F_1 \subseteq F_2$.
    \end{enumerate}
\end{proposition}

\begin{proof}
    \begin{enumerate}
        \item Let $\sigma \in \aut{K/ F_2}$. Let $a \in F_1$. Then $a \in F_2$ so 
        $$\sigma(a) = a$$
        Hence, $\sigma$ fixes $F_1$ and 
        $$\aut{K/F_2} \subseteq \aut{K/F_1}$$

        \item Let $\sigma \in H_1.$ Let $a \in F_1$. Since $H_1 \subseteq H_2$, $\sigma \in H_2$ so
        $$a = \sigma(a) \in F_2$$
    \end{enumerate}
\end{proof}

\begin{example}
    Let $K = \Q(\sqrt{2})$. Then the fixed field of 
    $$\aut{\Q(\sqrt{2})} = \aut{\Q(\sqrt{2})/\Q} = \set{1, \sigma}$$
    is the set of all $a+b\sqrt{2} \in \Q(\sqrt{2})$ such that 
    \begin{align*}
        1(a + b\sqrt{2}) &= a + b\sqrt{2} \\
        \sigma(a + b\sqrt{2}) &= a + b\sqrt{2}
    \end{align*}
    Thus, we must have that 
    \begin{align*}
        a-b\sqrt{2} &= a+b\sqrt{2} \\
        \implies -b\sqrt{2} &= b\sqrt{2} \\
        \implies -b &= b \\
        \implies b &= 0
    \end{align*}
    So, the fixed field of $\aut{\Q(\sqrt{2})/\Q}$ is just $\Q$.
\end{example}

\begin{example}
    Let $K = \Q(\sqrt{2})$. The fixed field of $\aut{\Q(\sqrt[3]{2})/\Q} = 1$ is $Q(\sqrt[3]{2})$. Notice that for $\Q(\sqrt{2})$ we have

    \begin{center}
        \begin{tabular}{ccc}
            % Header Row
            \textbf{Subgroups of $\aut{\Q({\sqrt{2}})/\Q}$:} & & \textbf{Fixed fields:} \\ 
            \noalign{\medskip} % Adds a bit more space after the header
            
            % Data Row
            $\set{1}$ & $\longleftrightarrow$ & $\Q(\sqrt{2})$ \\
            $\set{1, \sigma}$ & $\longleftrightarrow$ & $\Q$
        \end{tabular}
    \end{center}

    With $\Q(\sqrt[3]{2})$, we have 
    \begin{center}
        \begin{tabular}{ccc}
            % Header Row
            \textbf{Subgroups of $\aut{\Q({\sqrt{2}})/\Q}$:} & & \textbf{Fixed fields:} \\ 
            \noalign{\medskip} % Adds a bit more space after the header
            
            % Data Row
            $\set{1}$ & $\longleftrightarrow$ & $\Q(\sqrt[3]{2})$
        \end{tabular}
    \end{center}
    but we have no set of automorphisms which fix $\Q$.
\end{example}

\begin{proposition}
    Let $E$ be the splitting field over $F$ of some polynomial $f(x) \in F[x]$. Then
    $$|\aut{E/F}| \leq [E:F]$$
    with equality when $f(x)$ is separable over $F$.
\end{proposition}

\begin{definition}[Galois Extensions and Groups] \leavevmode \\
    Let $K/F$ be a finite extension. Then $K$ is said to be Galois over $F$ ($K/F$ is a Galois Extension) if $|\aut{K/F}| = [K:F]$. If $K/F$ is Galois, the group of automorphisms $\aut{K/F}$ is called the Galois group of $K/F$, denoted $\Gal{K/F}$.
\end{definition}

\begin{corollary}
    If $K$ is the splitting field over $F$ of a separable polynomial $f(x)$, then $K/F$ is Galois.
\end{corollary}

\begin{definition}[Galois Group of a Polynomial] \leavevmode \\
    If $f(x)$ is a separable polynomial over $F$, then the Galois group of $f(x)$ over $F$ is $\Gal{K/F}$ where $K$ is the splitting field of $f(x)$ over $F$.
\end{definition}

\begin{example}
    Show $\Q(\sqrt{2})/\Q$ is Galois. \\
    We know that $[\Q(\sqrt{2}):\Q] = 2$ and that $\aut{\Q(\sqrt{2})/\Q} = \set{1, \sigma}$. Hence, $\Q(\sqrt{2})/\Q$ is a Galois extension.

    Alternatively, we know that $\Q(\sqrt{2})$ is the splitting field of $x^2-2 \in \Q[x]$. Hence, $\Q(\sqrt{2})$ is Galois over $\Q$.
\end{example}

\begin{example}
    The extension $\Q(\sqrt[3]{2})/\Q$ is not Galois since
    $$[\Q(\sqrt[3]{2}):\Q] = 3 \text{ but } \aut{\Q(\sqrt[3]{2})/\Q} = 1$$
\end{example}

\begin{example}
    The extension $\Q(\sqrt{2}, \sqrt{3})$ is Galois since it is the splitting field of the separable polynomial $(x^2-2)(x^2-3) \in \Q [x]$. Calculating the Galois group:
    
    Any $\sigma \in \Gal{\Q(\sqrt{2},\sqrt{3})/\Q}$ can be determined by its action on $\sqrt{2}$ and $\sqrt{3}$. We know that $\sqrt{2}$ and $\sqrt{3}$ must map to $\pm \sqrt{2}$ and $\pm \sqrt{3}$, respectively. Hence, the only possible automorphisms are
    \begin{center}
        \begin{tabular}{cc}
            % Data Row
            $\begin{cases*}
                \sqrt{2} \mapsto \sqrt{2} \\
                \sqrt{3} \mapsto \sqrt{3}
            \end{cases*}$ & 
            $\begin{cases*}
                \sqrt{2} \mapsto -\sqrt{2} \\
                \sqrt{3} \mapsto \sqrt{3}
            \end{cases*}$ \\
            $\begin{cases*}
                \sqrt{2} \mapsto \sqrt{2} \\
                \sqrt{3} \mapsto -\sqrt{3}
            \end{cases*}$ &
            $\begin{cases*}
                \sqrt{2} \mapsto -\sqrt{2} \\
                \sqrt{3} \mapsto -\sqrt{3}
            \end{cases*}$
        \end{tabular}
    \end{center}
    Let
    \begin{align*}
        \sigma = 
        \begin{cases*}
            \sqrt{2} \mapsto -\sqrt{2} \\
            \sqrt{3} \mapsto \sqrt{3}
        \end{cases*}, ~~\tau = 
        \begin{cases*}
            \sqrt{2} \mapsto \sqrt{2} \\
            \sqrt{3} \mapsto \sqrt{3}
        \end{cases*}
    \end{align*}
    Notice
    $$\sigma^2(\sqrt{2}) = \sqrt{2} \text{ and } \sigma^2(\sqrt{3}) = \sqrt{3} \implies \sigma^2 = 1$$
    Similarly, $\tau^2 = 1$. Furthermore,
    $$\sigma(\tau(\sqrt{2})) = -\sqrt{2}, ~~\sigma(\tau(\sqrt{3}))=-\sqrt{3}$$
    so
    $$
    \sigma \tau = \begin{cases*}
        \sqrt{2} \mapsto -\sqrt{2} \\
        \sqrt{3} \mapsto -\sqrt{3}
    \end{cases*}
    $$
    Lastly, $\sigma \tau = 1$ so 
    $$\Gal{\Q(\sqrt{2}, \sqrt{3})/\Q} \cong \Z_2 \times \Z_2$$
\end{example}

\begin{example} Let's look at the corresponding fixed fields of the previous example.

    \begin{center}
        \begin{tabular}{ccc}
            \textbf{Subgroups of $\Gal{\Q({\sqrt{2}}, \sqrt{3})/\Q}$:} & & \textbf{Corresponding fixed fields:} \\ 
            \noalign{\medskip}
            % Data Row
            $1$ & $\longleftrightarrow$ & $\Q(\sqrt{2}, \sqrt{3})$ \\
            $\cyc{\sigma}$ & $\longleftrightarrow$ & $\Q(\sqrt{3})$ \\
            $\cyc{\tau}$ & $\longleftrightarrow$ & $\Q(\sqrt{2})$ \\
            $\cyc{\sigma \tau}$ & $\longleftrightarrow$ & $\Q(\sqrt{6})$ \\ 
            $\set{1, \sigma, \tau, \sigma \tau}$ & $\longleftrightarrow$ & $\Q$
        \end{tabular}
    \end{center}
    To find the fixed field for $\cyc{\sigma}$, note that
    \begin{align*}
        \sigma(a+b\sqrt{2}+c\sqrt{3}+d\sqrt{6}) &= a + b\sigma(\sqrt{2})+c\sigma(\sqrt{3})+d\sigma(\sqrt{6}) \\
        &= a-b\sqrt{2}+c\sqrt{3}+d\sigma(\sqrt{2})\sigma(\sqrt{3}) \\
        &= a-b\sqrt{2}+c\sqrt{3}-d\sqrt{6}
    \end{align*}
    An element in $\Q(\sqrt{2}, \sqrt{3})$ is fixed by $\sigma$ if 
    \begin{align*}
        a+b\sqrt{2} + c\sqrt{3} + d\sqrt{6} &= a-b\sqrt{2}+c\sqrt{3}-d\sqrt{6} \\
        \implies b=-b &\text{ and } d = -d \\
        \implies b=&d=0
    \end{align*}
    so $\sigma(a+b\sqrt{2}+c\sqrt{3}+d\sqrt{6}) = a+c\sqrt{3} \in \Q(\sqrt{3})$.

    Similarly, for $\cyc{\tau},$
    \begin{align*}
        \tau(a+b\sqrt{2}+c\sqrt{3}+d\sqrt{6}) &= a + b\sqrt{2}-c\sqrt{3}-d\sqrt{6} \\
        \implies c = &0 = d \\
        \implies \tau(a+b\sqrt{2}+c\sqrt{3}+d\sqrt{6})&= a + b\sqrt{2} \in \Q(\sqrt{2})
    \end{align*}
\end{example}

\begin{example}
    Let $p(x) = x^3-2 \in \Q[x]$. We have found the splitting field of $p(x)$ over $\mathbb{Q}$ to be $\Q(\sqrt[3]{2}, \omega)$ where $\omega = \frac{-1+\sqrt{3}i}{2}.$

    Notice that $x^3-2$ has three distinct roots, so $p(x)$ is separable and so its splitting field is Galois over $\Q$. We also know
    $$[\Q(\sqrt[3]{2}, \omega):\Q] = 6 = \left| \Gal{\Q(\sqrt[2]{3}, \omega)/\Q}\right|.$$
    Suppose we write the splitting field as
    $$\Q(\sqrt[3]{2}, \sqrt[3]{2}\omega)$$
    We know that any automorphism will send roots of minimal polynomials to other roots. Written this way, we may have 
    \[
        \begin{tikzcd}[column sep=large]
            & \sqrt[3]{2} \\
            \sqrt[3]{2} \arrow[ru, mapsto] \arrow[r, mapsto] \arrow[rd, mapsto] & \sqrt[3]{2} \omega \\
            & \sqrt[3]{2} \omega^2
        \end{tikzcd}, 
        \begin{tikzcd}[column sep=large]
            & \sqrt[3]{2} \\
            \sqrt[3]{2}\omega \arrow[ru, mapsto] \arrow[r, mapsto] \arrow[rd, mapsto] & \sqrt[3]{2} \omega \\
            & \sqrt[3]{2} \omega^2
        \end{tikzcd}
    \]

    This gives $3 \cdot 3 = 9$ possibilities for automorphisms, which is larger than the order of the Galois group. 

    Using $\Q(\sqrt[3]{2}, \omega)$, we have a more compatible number of potential automorphisms. Any automorphism can send
        \[
        \begin{tikzcd}[column sep=large]
            \omega \arrow[r, mapsto] \arrow[rd, mapsto] & \omega \\
             & \omega^2
        \end{tikzcd}, 
        \begin{tikzcd}[column sep=large]
             & \sqrt[3]{2} \\
            \sqrt[3]{2} \arrow[ru, mapsto] \arrow[r, mapsto] \arrow[rd, mapsto] & \sqrt[3]{2} \omega \\
            & \sqrt[3]{2} \omega^2
        \end{tikzcd}
    \]
    ($m_{\omega, \Q}(x) = x^2 + x + 1$, $m_{\sqrt[3]{2}, \Q}(x) = x^3-2$).
    This gives six automorphisms, which is exactly the size of the Galois group. Let
    \begin{align*}
        1 &: \begin{cases*}
            \sqrt[3]{2} \mapsto \sqrt[3]{2} \\
            \omega \mapsto \omega
        \end{cases*} \\
        \sigma &: \begin{cases*}
            \sqrt[3]{2} \mapsto \sqrt[3]{2}\omega \\
            \omega \mapsto \omega
        \end{cases*} &&
        \tau: \begin{cases*}
            \sqrt[3]{2} \mapsto \sqrt[3]{2} \\
            \omega \mapsto \omega^2 = -1-\omega
        \end{cases*}
    \end{align*}
    \begin{align*}
                \sigma^2 &: \begin{cases*}
            \sqrt[3]{2} \mapsto \sigma(\sqrt[3]{2}\omega) \mapsto \sqrt[3]{2}\omega^2 = \sqrt[3]{2}(-1-\omega) \\
            \omega \mapsto \omega
        \end{cases*} &&
        \tau^2: \begin{cases*}
            \sqrt[3]{2} \mapsto \sqrt[3]{2} \\
            \omega \mapsto \left(\omega^2\right)^2 = \omega^4 = \omega
        \end{cases*} \\
        \sigma^3 &: \begin{cases*}
            \sqrt[3]{2} \mapsto \sqrt[3]{2}\omega(-1-\omega) = \sqrt[3]{2}(-\omega + 1 + \omega) = \sqrt[3]{2} \\
            \omega \mapsto \omega
        \end{cases*}
    \end{align*}
    Notice that $\sigma^3 = 1 = \tau^2$, so $|\sigma| = |\cyc{\sigma}| = 3$ and $|\tau| = |\cyc{\tau}| = 2$.

    \begin{align*}
        \sigma \tau &: \begin{cases*}
            \sqrt[3]{2} \mapsto \sqrt[3]{2} \omega \\
            \omega \mapsto -1-\omega
        \end{cases*} &&
        \tau \sigma: \begin{cases*}
            \sqrt[3]{2} \mapsto \sqrt[3]{2}(-1-\omega) \\
            \omega \mapsto (-1-\omega)
        \end{cases*}
    \end{align*}
    $\omega \tau \neq \tau \omega,$ so the Galois group is not abelian.
    \begin{align*}
        \tau\sigma^2 &: \begin{cases*}
            \sqrt[3]{2} \mapsto \sqrt[3]{2}\omega \\
            \omega \mapsto -1-\omega
        \end{cases*}
    \end{align*}
    $$\sigma\tau = \tau \sigma^2 \iff \sigma \tau = \tau \sigma^{-1}$$
    Hence,
    $$\Gal{\Q(\sqrt[3]{2}, \omega)/\Q} \cong \cyc{\sigma, \tau : \sigma^3 = \tau^2 = 1, \sigma \tau = \tau \sigma^{-1}} \cong D_6$$
\end{example}

\begin{example}
    Let's compute the fixed field of $\cyc{\sigma}$ from the previous example.

    We need to see what is fixed by the generator $\sigma$. Lets look at the basis for $\Q(\sqrt[3]{2}, \omega)$:
    $$\set{1, \sqrt[3]{2}, \left(\sqrt[3]{2}\right)^2, \omega, \sqrt[3]{2}\omega, \left(\sqrt[3]{2}\right)^2\omega}.$$
    Then,
    \begin{align*}
        \sigma(1) &= 1 \\
        \sigma\left(\sqrt[3]{2}\right) &= \sqrt[3]{2}\omega \\
        \sigma\left((\sqrt[3]{2}) ^2\right) &= \sigma\left(\sqrt[3]{2}\right)\sigma\left(\sqrt[3]{2}\right) = \left(\sqrt[3]{2}\omega\right)\left(\sqrt[3]{2}\omega\right) = \sqrt[3]{4}(-1-\omega) \\
        \sigma(\omega) &= \omega \\
        \sigma\left(\sqrt[3]{2} \omega\right) &= \sigma\left(\sqrt[3]{2}\right)\sigma(\omega) = \left(\sqrt[3]{2} \omega\right)\left(\omega\right) = \sqrt[3]{2}(-1-\omega) \\
        \sigma\left((\sqrt[3]{2})^2 \omega\right) &= \sigma\left((\sqrt[3]{2})^2\right)\sigma(\omega) = \left(\sqrt[3]{4}(-1-\omega)\right)(\omega) = \sqrt[3]{4}(-\omega + 1 + \omega) = \sqrt[3]{4}
    \end{align*}
    So,
    \begin{align*}
        \sigma\left(a+b\sqrt[3]{2} + c(\sqrt[3]{2})^2+d\omega + e\sqrt[3]{2} \omega + f (\sqrt[3]{2})^2 \omega\right)
        &= a + b \sqrt[3]{2} \omega + c \left((\sqrt[3]{2})^2\omega^2\right) + dw + e\sqrt[3]{2} \omega^2 + f (\sqrt[3]{2})^2 \\
        &= a + b \sqrt[3]{2} \omega - c(\sqrt[3]{2})^2(1+\omega) + d\omega - e \sqrt[3]{2}(1+w) + f(\sqrt[3]{2})^2 \\
        &= a - e \sqrt[3]{2} + (f-c)(\sqrt[3]{2})^2 +d\omega + (b-e)\sqrt[3]{2} \omega -c (\sqrt[3]{2})^2 \omega
    \end{align*}
    For our element to be fixed, we need
    \begin{align*}
        &a = a &&f=-c \\
        &b=-e &&c=f-c \\
        &d=d &&e=b-e
    \end{align*}
    We have
    $$e=b-e = -e-e \implies e = 0 \implies b = 0$$
    Similarly,
    $$c=0 \implies f = 0$$
    Hence, any elements fixed by $\sigma$ are of the form $a + dw$. The fixed field of $\cyc{\sigma}$ is $\Q(\omega)$.
\end{example}

\begin{example}
    $\Q(\sqrt{2}, \sqrt[3]{3}, \sqrt[5]{5},...) = \bigcap \limits_{\substack{F \supseteq \Q \\ \sqrt[p_i]{p_i} \in F}} F$
    Note that
    $$\bigcap_{\substack{F \supseteq \Q \\ \sqrt[p_i]{p_i} \in F}} F = \bigcup_{i=1}^\infty K_i \text{ where } K_i = \Q\left(\sqrt[2]{2}, \sqrt[3]{3}, \sqrt[5]{5},...,\sqrt[p_i]{p_i}\right)$$
\end{example}